\section{Related Works}
% \textbf{Reconstruction-based Driving Simulators.} 
% In recent years, 3D scene reconstruction based on Neural Radiance Fields (NeRF)\cite{mildenhall2020nerf} and 3D Gaussian Splatting (3DGS)\cite{kerbl20233d} have been fully developed. NeRF represents scenes as continuous volumetric radiance fields parameterized by multilayer perceptions. 3DGS represents scenes with clouds of anisotropic Gaussians, enabling rapid rasterization and real‐time rendering while maintaining high image rendering quality. Recently, many works have extended NeRF and 3DGS to unbounded, dynamic driving scenes\cite{chen2023periodic,huang2024s3gaussian,wu2023mars,guo2023streetsurf,turki2023suds,yang2023unisim,tonderski2024neurad,zhou2024drivinggaussian,yan2024street}. HUGS\cite{zhou2024hugs} incorporated inter-frame pixel optical flow and 3D semantic information into the scene reconstruction pipeline. OmniRe\cite{chen2025omnire} further introduced the non-rigid modeling of pedestrians and cyclists. However, due to the lack of observations of the 3D scene outside the original trajectory, reconstruction-based methods cannot render high-quality novel camera views from off-trajectory viewpoints.

% \textbf{Generation-reconstruction hybrid Driving Simulators.} 
% To overcome the constraints imposed by the original observation viewpoint distribution on the rendering range of NVS in driving scenarios, incorporating generative models into the NVS task for driving scenes has become a research hotspot. FreeSim\cite{fan2024freesim} adopts a progressive reconstruction strategy, which progressively adds generated images of unrecorded views into the reconstruction process, starting from slightly off-trajectory viewpoints and moving progressively farther away. Similarly, DriveDreamer4D\cite{zhao2024drive} and ReconDreamer\cite{Ni2024ReconDreamerCW} also incorporate the output of generative models as additional supervision for scene reconstruction, aiming to combine generative and reconstruction models. However, generative models often struggle to maintain consistent 3D scene details across different viewpoints, leading to under-converged 3D reconstruction and the blurring effect of scene details. Moreover, since these methods still adopt 3DGS as the final scene representation, it remains difficult for those methods to flexibly edit scene content, similar to that of pure reconstruction-based methods.

% \textbf{Diffusion based Driving Simulators.} 

 
% Some video generation methods like Cosmos~\cite {agarwal2025cosmos} and MaskGWM~\cite{ni2025maskgwm} can generate multi-view consistent videos given text prompts. However these method do not construct an explicit 4D representation of the scene, which means we can not control the camera pose explicitly. FreeVS\cite{wang2024freevs} utilizes the LiDAR point clouds as a prior to generative videos. Although we can control the target camera pose explicitly by controlling the lidar projection matrix, FreeVS is unable to recover information in the 3D scene that is not captured by the LiDAR point clouds, such as distant or elevated 3D content.

\textbf{Reconstruction-Based Driving Simulators.}  
Recent advancements in 3D scene reconstruction have been driven by Neural Radiance Fields (NeRF)~\cite{mildenhall2020nerf} and 3D Gaussian Splatting (3DGS)~\cite{kerbl20233d}. NeRF models scenes as continuous volumetric radiance fields, while 3DGS represents scenes using anisotropic Gaussian clouds. Building upon these foundations, numerous works have extended NeRF and 3DGS to unbounded, dynamic driving environments~\cite{chen2023periodic,huang2024s3gaussian,wu2023mars,guo2023streetsurf,turki2023suds,yang2023unisim,tonderski2024neurad,zhou2024drivinggaussian,yan2024street}. For example, OmniRe~\cite{chen2025omnire} introduces non-rigid modeling for dynamic actors such as pedestrians and cyclists. However, these reconstruction-based approaches are inherently limited by the original camera trajectory. Since they lack observations beyond the recorded viewpoints, they struggle to render high-quality novel views from significantly off-trajectory positions.

\textbf{Generation-Reconstruction Hybrid Driving Simulators.}  
To overcome the limitations imposed by the original viewpoint distribution, recent approaches have explored hybrid models that combine generative techniques with scene reconstruction for novel view synthesis (NVS). FreeSim~\cite{fan2024freesim} progressively incorporates generated views into reconstruction, beginning near the recorded trajectory and expanding outward. Similarly, DriveDreamer4D~\cite{zhao2024drive} and ReconDreamer~\cite{Ni2024ReconDreamerCW} use generative outputs as supervision signals to improve reconstruction quality. While these methods aim to blend the strengths of generative and reconstruction-based models, they still face challenges. Generative models often struggle to preserve consistent 3D structure across viewpoints, leading to under-converged geometry and blurred scene details. Moreover, as these methods ultimately rely on 3DGS as the scene representation, they inherit the same limitations in editability and flexibility as pure reconstruction-based approaches.

\textbf{Diffusion-Based Driving Simulators.}  
Diffusion-based generative models have also been explored for driving simulation. For instance, Vista~\cite{gao2024vista} and MaskGWM~\cite{ni2025maskgwm} generate driving videos from text prompts. However, these models do not construct an explicit 4D representation of the scene, limiting the ability to directly control camera poses. FreeVS~\cite{wang2024freevs} addresses this by incorporating LiDAR point clouds as priors to guide video generation, allowing explicit control of the target camera pose via the LiDAR projection matrix. Nevertheless, FreeVS cannot recover 3D content not captured by the LiDAR sensor, such as distant or elevated structures, which constrains the completeness of the generated scenes.